\scanpage{48}
% archive sheet number: 47
\begingroup\csname tagsleft@true\endcsname
loc. isomorphe à \(\{0,\infty,1,-1,i,-i\}\) de \(\mathbb P^1_S\), ou ce qui revient au \(=\), qui soit loc. isom. à la configuration octaédrale associée à une op. \inserted{\unsure{monom.}} d'un \(\underline V\) sur \(X\). Nous allons voir maintenant que cette donnée est essentiellement unique.

Pour ceci, pour un diviseur octaédral \(\deleted{\unsure{\Delta}}\)
\[
\Delta\subset X
\]
considérons le schéma en groupes
\[
\deleted{\unsure{G}}\Gamma=\underline{\operatorname{Aut}}_{\Delta}(X,\Delta), \tag*{(7)}
\]
Il est immédiat (du fait que \(\Delta\) est de degré \(\geq3\)) que \(\Gamma\) est \uline{net} sur \(S\), \inserted{car \(\subset\underline{\mathfrak S}_{\Delta}=\underline{\operatorname{Aut}}_S(\Delta)\)} Je dis que \(\Gamma\) est étale sur \(S\), et loc. isom au schéma en groupes \((\mathfrak S_4)_S\). Plus précisément, je vais définir sur \(\Delta\) \inserted{(comme schéma des sommets)} une structure d'o\deleted{\unsure{ct}} schéma en octaèdres combinatoires; de telle façon que l'on ait
\[
\begin{aligned}
\Gamma&=\underline{\operatorname{Aut}}^{+}_{\mathrm{oct.}}(\Delta)
\qquad\bigl(\text{en tant que}\\
&&\text{sous-schémas de}\\
&&\underline{\operatorname{Aut}}_S(\Delta)\bigr)
\end{aligned} \tag*{(8)}
\]
Pour ceci, on \deleted{\unsure{va décrire une}} \inserted{note qu'une} structure octaédrale sur un ens. à six éléments \(\Delta\) est bien déterminée quand on connaît le groupe de ses automorphismes \inserted{directs} (en tant que \unsure{sous-groupes|sous-groupe} de \(\mathfrak S_{\Delta}\)), car
\endgroup
